% Part: proof-theory
% Chapter: cut-elimination
% Section: introduction

\documentclass[../../../include/open-logic-section]{subfiles}

\begin{document}

\olfileid{pt}{cut}{int}

\olsection{تمهيد في حذف القطع}

قاعدة~\Cut{} هي
\[
\Axiom$\Gamma \fCenter \Delta, !A$
\Axiom$!A, \Pi \fCenter \Lambda$
\RightLabel{\Cut}
\BinaryInf$\Gamma, \Pi \fCenter \Delta, \Lambda$
\DisplayProof
\]
(وتسمى~!!{formula}~$!A$ \emph{صيغة القطع}.)

أدخل غنتسن قاعدة~\Cut{} في حساب المتتاليات الأصلي~\Log{LK}. ومؤدى
\emph{مبرهنته الرئيسة} (\emph{Hauptsatz}) أن كل متتالية~!!{provable}
في~\Log{LK} تبقى~!!{provable} من غير استعمال قاعدة~\Cut{}؛ أي إن
\Cut{} «تحذف» من براهين~\Log{LK}. وأما النظامان~\Log{G1c} و\Log{G3c}
فليس فيهما قاعدة~\Cut، غير أن السؤال بعينه يصح فيهما: هل تحذف~\Cut{}
من~!!{proof}s التي تضمها إلى سائر قواعد~\Log{G1c} (أو قواعد~\Log{G3c})؟
وهذا، بعبارتنا الاصطلاحية، هو السؤال: هل قاعدة~\Cut{} مقبولة في
\Log{G1c} (أو~\Log{G3c})؟

ومبرهنة حذف القطع من أجل نتائج نظرية البرهان وأشهرها. وأول ما تفيده أن
\Log{G1c} و\Log{G3c} مستغنيان عن قاعدة يخيل، في أول النظر، أنها لا غنى
عنها. وسبب ذلك أن صورنة البراهين الجارية تحتاج إلى حساب استعمال اللمم؛
فالرياضي يثبت نتيجة ثم يبني عليها نتيجة أخرى. فقد نصوغ أولًا برهان
النتيجة~$L$، وهي اللمّة، في~!!{proof} للمتتالية~$\Sequent L$، ثم نصوغ
برهان النتيجة الثانية~$R$، مع استعمال~$L$، في~!!{proof} للمتتالية
$L \Sequent R$. فبأي وجه نثبت أن لـ$R$ برهانًا، أي أن للمتتالية
$\Sequent R$~!!a{proof}؟ لا بد من وصل~!!{proof}s الاثنين، وهذا هو ما
تتيحه قاعدة~\Cut{}.

قد تقدم أن~\Log{G1c} و\Log{G3c} تامان، أي إنهما يبرهنان كل متتالية
صادقة، ويمكن إقامة هذا مباشرة. لكن التمام في عصر غنتسن لم يكن معلومًا
إلا لأنظمة الاشتقاق المسلّمية. ولذلك أثبت غنتسن تمام~\Log{LK} بأن ترجم
!!{derivation}s النظام المسلّمي إلى~!!{proof}s في~\Log{LK}، وكانت
\Cut{} ركنًا في تلك الترجمة. ثم إن شأن حذف القطع لا يختص بالمنطق
الكلاسيكي؛ فإن حسابات المتتاليات موجودة لمنطقات كثيرة، وقد يعسر أو يتعذر
في بعضها البرهان المباشر على تمام الحساب. فلا يبقى لإثبات التمام إلا
الترجمة من أنظمة أخرى باستعمال~\Cut. وعند ذلك ينبغي برهنة حذف القطع
برهانيًا، ليظهر أن~\Cut{} فضلة وأن النظام الخالي منها تام.

ولحذف القطع شأن آخر، إذ تترتب عليه نتائج مستقلة الأهمية. وسنعرض منها
لمّة كريغ للاستيفاء ومبرهنة هيربراند.

وعمدة البراهين المعتادة لحذف القطع أن تبين كيف يحول~!!a{proof} يستعمل
\Cut{} إلى برهان مستغن عنها. وأكثر هذه التحويلات «موضعي»: نأخذ قطعة من
!!a{proof}، هي في الغالب~!!{proof} جزئي آخره استدلال~\Cut{}، ونقيم
مكانها~!!{proof} جزئيًا للمتتالية نفسها، حذف منه ذلك الاستدلال أو عوض
باستدلالات غيره. ومن هذه الحجج تنشأ خوارزميات تحول، خطوة فخطوة،
!!{proof}s المشتملة على~\Cut{} إلى براهين خالية منها؛ ولذلك تفيد معلومات
زائدة عند تطبيق إجراءات حذف القطع. فمثلًا قد يكون البرهان النموذجي
النظري للمّة الاستيفاء على الصورة: «إذا كانت~$!A \lif !B$ صادقة، وجدت
صيغة استيفاء» (وليس المقصود الآن بيان معناها). أما الحجة البرهانية
المستعملة لحذف القطع فتعطي خوارزمية مدخلها~!!a{proof} لـ$!A \lif !B$
ومخرجها صيغة استيفاء. فهي، بخلاف الحجة النموذجية النظرية،
\emph{بنائية}.

ولبرهان المبرهنة طريقان مشهوران. أما طريق غنتسن فيبين إمكان إزالة
القطوع العليا من البرهان، ثم يزيلها متعاقبة حتى تنفد. وأما الطريق الذي
أصله لشوتّه وتايت فيقصد أعقد القطوع في~!!a{proof}، فيزيل تباعًا القطوع
القصوى، أي التي لا يوجد قطع ذو~!!{formula} أعمق منها. ولكل طريق موضع قوة
وموضع ضعف؛ فقد يتيسر أحدهما في نظام ويتعسر الآخر أو يمتنع. ولهذا نصفهما
معًا، ونبرهن حذف القطع لـ\Log{G3c}. على أننا لا نبين حذف~\Cut{} نفسها،
بل حذف
\emph{القطع المشترك السياق}~\CutCS:
\[
\Axiom$\Gamma \fCenter \Delta, !A$
\Axiom$!A, \Gamma \fCenter \Delta$
\RightLabel{\CutCS}
\BinaryInf$\Gamma \fCenter \Delta$
\DisplayProof
\]
ومتى أجاز حساب المتتاليات الإضعاف والتقليص، سواء أكانا قاعدتين فعليتين
كما في~\Log{G1c} أم قاعدتين مقبولتين كما في~\Log{G3c}، حاكت~\Cut{}
قاعدة~\CutCS وحاكتها الثانية. وإنما نؤثر~\CutCS{} على~\Cut لأن وصف
الطريقة الأولى في حذف القطع أيسر معها، ولأن الطريقة الثانية لا تصح إلا
بها.

\end{document}
